\documentclass[10pt, a4paper]{article}
\usepackage{graphicx} 
\usepackage{listings}
\usepackage{xcolor} 
\usepackage[table]{xcolor} 
\usepackage{amsmath}
\usepackage{amssymb}
\usepackage{wrapfig}
\usepackage{parskip}
\usepackage{float}
\usepackage{hyperref}
\usepackage[utf8]{inputenc}
\usepackage[T1]{fontenc}
\usepackage{siunitx}
\sisetup{output-decimal-marker = {,}}
\usepackage[swedish]{babel}
\usepackage{array}
\usepackage{adjustbox}
\usepackage[margin=2.5 cm]{geometry}
\usepackage{appendix}
\usepackage{csquotes}
\usepackage{comment}
\usepackage[
    sorting=none
]{biblatex}
\usepackage{titling}

\addbibresource{../references.bib}

\lstdefinestyle{mypython}{
    language=Python,
    basicstyle=\ttfamily\small,
    backgroundcolor=\color{gray!10},
    keywordstyle=\color{blue}\bfseries,
    %stringstyle=\color{orange!85!black},
    commentstyle=\color{green!50!black},
    numberstyle=\tiny,
    numbers=left,
    stepnumber=1,
    numbersep=8pt,
    showstringspaces=false,
    frame=single,
    breaklines=true,
    breakatwhitespace=true,
    tabsize=4,
    morekeywords={self},  % Lägg till fler Python-nyckelord om nödvändigt
    literate={ö}{{\"o}}1 {å}{{\aa}}1 {ä}{{\"a}}1 {\\}{{\textbackslash}}1,
    }
    
\begin{document}

\pagenumbering{arabic}
\thispagestyle{empty}
\begin{center}
{\huge\bf Laborationsrapport - Kvantfysik:\\}
{\Large\bf Studier av atomära spektrum
}
\end{center}
\vspace{1.0cm}
\begin{center}
Julius Berg (juliusbe), Benjamin Johansson (benjjoh), \par
Program: Teknisk fysik. \par
Kurs: TIF097 - Experimentell fysik 2.
\end{center}
\vspace{1.0cm}
\centerline{\bf Sammandrag}
 
\newpage
\tableofcontents
\newpage

\section{Inledning}

Alla grundämnen interagerar på olika sätt med elektromagnetisk strålning.
Detta sker både via absorption, då en elektron absorberar en foton och ökar i energi, eller emission, 
då en elektron minskar i energi och emitterar en foton med motsvarande energi\cite{rana_quantum_electronics}.
Denna växelverkan sker dock endast vid specifika frekvenser, då elektroner i en atom endast kan existera
på specifika energinivåer.
När en elektron övergår från en energinivå till en energinivå med lägre energi frigörs därmed en foton som har en våglängd som ges av energiskillnaden mellan nivåerna.
Detta ger upphov till en uppsättning fotoner med en bestämd våglängd i ett spektrum, så kallade spektrallinjer\cite{TKACHENKO2006107}.

Eftersom olika grundämnen har olika elektronstrukturer har varje ämne en unik uppsättning möjliga energiövergångar.
Detta innebär att specifika grundämnen producerar unika uppsättningar av spektrallinjer.
Denna unika uppsättning fungerar som ett fingeravtryck för det specifika grundämnet, och detta faktum är mycket viktigt inom till exempel astronomi, 
där emissionsspektrum används för att bestämma vilka grundämnen och molekyler som finns i exoplaneters atmosfärer\cite{tessenyi_molecular_detectability}.

Syftet med denna studie är att för ett bestämt grundämne bestämma emissionsspektrum och utifrån detta emissionsspektrum konstruera grundämnets energinivådiagram.
\section{Bakgrund}

\section{Metod}

\subsection{Försöksuppställning}

\subsection{Genomförande}

\section{Förväntade resultat}

\section{Diskussion}

\newpage

\printbibliography

\end{document}